\documentclass[10pt,a4paper]{amsart}
\usepackage{amsmath,amssymb,amsthm,graphicx}
\usepackage{verbatim}
\newenvironment{reference}{\comment}{\endcomment}
%\newenvironment{reference}{}{}
\newenvironment{slogan}{\comment}{\endcomment}
\newenvironment{history}{\comment}{\endcomment}

% For commutative diagrams we use Xy-pic
\usepackage[all]{xy}

% We use 2cell for 2-commutative diagrams.
\xyoption{2cell}
% Theorem environments.
%
\theoremstyle{plain}
\newtheorem{theorem}[subsection]{Theorem}
\newtheorem{proposition}[subsection]{Proposition}
\newtheorem{lemma}[subsection]{Lemma}

\theoremstyle{definition}
\newtheorem{definition}[subsection]{Definition}
\newtheorem{example}[subsection]{Example}
\newtheorem{exercise}[subsection]{Exercise}
\newtheorem{situation}[subsection]{Situation}

\theoremstyle{remark}
\newtheorem{remark}[subsection]{Remark}
\newtheorem{remarks}[subsection]{Remarks}

\numberwithin{equation}{subsection}

% Macros
%
\def\lim{\mathop{\mathrm{lim}}\nolimits}
\def\colim{\mathop{\mathrm{colim}}\nolimits}
\def\Spec{\mathop{\mathrm{Spec}}}
\def\Hom{\mathop{\mathrm{Hom}}\nolimits}
\def\Ext{\mathop{\mathrm{Ext}}\nolimits}
\def\SheafHom{\mathop{\mathcal{H}\!\mathit{om}}\nolimits}
\def\SheafExt{\mathop{\mathcal{E}\!\mathit{xt}}\nolimits}
\def\Sch{\mathit{Sch}}
\def\Mor{\mathop{\mathrm{Mor}}\nolimits}
\def\Ob{\mathop{\mathrm{Ob}}\nolimits}
\def\Sh{\mathop{\mathit{Sh}}\nolimits}
\def\NL{\mathop{N\!L}\nolimits}
\def\CH{\mathop{\mathrm{CH}}\nolimits}
\def\proetale{{pro\text{-}\acute{e}tale}}
\def\etale{{\acute{e}tale}}
\def\QCoh{\mathit{QCoh}}
\def\Ker{\mathop{\mathrm{Ker}}}
\def\Im{\mathop{\mathrm{Im}}}
\def\Coker{\mathop{\mathrm{Coker}}}
\def\Coim{\mathop{\mathrm{Coim}}}

% Boxtimes
%
\DeclareMathSymbol{\boxtimes}{\mathbin}{AMSa}{"02}

%
% Macros for moduli stacks/spaces
%
\def\QCohstack{\mathcal{QC}\!\mathit{oh}}
\def\Cohstack{\mathcal{C}\!\mathit{oh}}
\def\Spacesstack{\mathcal{S}\!\mathit{paces}}
\def\Quotfunctor{\mathrm{Quot}}
\def\Hilbfunctor{\mathrm{Hilb}}
\def\Curvesstack{\mathcal{C}\!\mathit{urves}}
\def\Polarizedstack{\mathcal{P}\!\mathit{olarized}}
\def\Complexesstack{\mathcal{C}\!\mathit{omplexes}}
% \Pic is the operator that assigns to X its picard group, usage \Pic(X)
% \Picardstack_{X/B} denotes the Picard stack of X over B
% \Picardfunctor_{X/B} denotes the Picard functor of X over B
\def\Pic{\mathop{\mathrm{Pic}}\nolimits}
\def\Picardstack{\mathcal{P}\!\mathit{ic}}
\def\Picardfunctor{\mathrm{Pic}}
\def\Deformationcategory{\mathcal{D}\!\mathit{ef}}
\usepackage{cleveref}
\crefname{originalreference}{original reference}{original references}
\Crefname{originalreference}{Original reference}{Original references}
\makeatletter
\newlabel{definition-limit-spectral-sequence}{{011O}{source}{}{sourcepointer-1}{}}
\newlabel{definition-limit-spectral-sequence@cref}{{[originalreference][0][]011O}{[1][0][]source}{}{}{}}
\makeatother
\crefname{originalequation}{original equation}{original equations}
\Crefname{originalequation}{Original equation}{Original equations}
\makeatletter
\newlabel{equation-graded-cohomology}{{0BDT}{source}{}{sourcepointer-2}{}}
\newlabel{equation-graded-cohomology@cref}{{[originalequation][0][]0BDT}{[1][0][]source}{}{}{}}
\makeatother
\crefname{originalreference}{original reference}{original references}
\Crefname{originalreference}{Original reference}{Original references}
\makeatletter
\newlabel{lemma-spectral-sequence-filtered-complex}{{012M}{source}{}{sourcepointer-3}{}}
\newlabel{lemma-spectral-sequence-filtered-complex@cref}{{[originalreference][0][]012M}{[1][0][]source}{}{}{}}
\makeatother
\title{The limit term of the spectral sequence of a filtered complex and the graded cohomology}
\author{The Stacks Project Authors}
\date{stacks-project:homology.tex:local-snapshot-20261012; GNU FDL 1.2 or later}
\begin{document}
\maketitle
\noindent\textit{Editorial note.} All delivered mathematical text is the author's own, copied from the pinned Stacks Project chapter file homology.tex; no statement or proof step was rewritten, reordered or removed.\par\medskip
\noindent\textit{Editorial note.} The counted proof is the author's complete argument: the existence of the bidegree parts of the limit cycle and boundary objects, the explicit quotient description of the limit term, and the two inclusions comparing it with the graded cohomology.\par\medskip
\noindent\textit{Editorial note.} The label lemma-compute-cohomology-filtered-complex names the counted statement, and the proof environment immediately following it in the source is delivered verbatim as its proof.\par\medskip
\noindent\textit{Editorial note.} The counted statement is carried by the author's own lemma environment in the chapter source.\par\medskip
\noindent\textit{Editorial note.} Pool: stacks-homology, secondary mining (previous owner lane-s4 kept 63 cards; cross-line claim\_key scan done, its DISPOSITIONS-stacks-homology.json aligned). Pre-open gate: python operations/agent-authoring/tools/precard\_gate.py --source homology.tex --label lemma-compute-cohomology-filtered-complex --start 6404 --end 6476 --local-id 28179 returned RESULT=OK (occupancy FREE; the gate additionally reported a published-overlap note for lines 6463-6542, which lies beyond this item's own statement and is left untouched).\par\medskip
\noindent\textit{Editorial note.} The statement and proof are copied from the pinned chapter file; the cited inputs are kept as external pointers with their canonical Stacks tags.\par\medskip
\noindent\textit{Allowed standard prerequisites (editorial).}\begin{itemize}
\item Construction of the spectral sequence of a filtered complex and of its limit term. Same chapter, the lemma constructing the spectral sequence of a filtered complex and the definition of the limit of a spectral sequence. Provides the bigraded objects, the pages and the definition of the limit term whose bidegree parts are computed here.
\item The graded pieces of the cohomology of a filtered complex. Same chapter, the identification of the graded pieces of the cohomology of a filtered complex as a quotient of a subquotient of the cycle and boundary objects. Supplies the equality with which the subquotient computation is compared at the end of the proof.
\end{itemize}\medskip
\ifdefined\hypertarget\hypertarget{sourcepointer-1}{}\fi
\noindent\textit{Source reference.} Original reference 011O: sources/stacks-homology/pinned-original-source.tex, source lines 5436--5436. Prerequisite pointer invoked at proof line(s) 6421; it names a standard fact of the same chapter that the counted argument uses, and every step of the counted proof built on it is retained verbatim in the delivered body.\par\medskip
\ifdefined\hypertarget\hypertarget{sourcepointer-2}{}\fi
\noindent\textit{Source reference.} Original equation 0BDT: sources/stacks-homology/pinned-original-source.tex, source lines 6398--6398. Prerequisite pointer invoked at proof line(s) 6475; it names a standard fact of the same chapter that the counted argument uses, and every step of the counted proof built on it is retained verbatim in the delivered body.\par\medskip
\ifdefined\hypertarget\hypertarget{sourcepointer-3}{}\fi
\noindent\textit{Source reference.} Original reference 012M: sources/stacks-homology/pinned-original-source.tex, source lines 6266--6266. Prerequisite pointer invoked at proof line(s) 6423; it names a standard fact of the same chapter that the counted argument uses, and every step of the counted proof built on it is retained verbatim in the delivered body.\par\medskip
\begin{lemma}
\label{lemma-compute-cohomology-filtered-complex}
Let $\mathcal{A}$ be an abelian category. Let $(K^\bullet, F)$ be a filtered
complex of $\mathcal{A}$. If $Z_\infty^{p, q}$ and $B_\infty^{p, q}$ exist
(see proof), then
\begin{enumerate}
\item the limit $E_\infty$ exists and is a bigraded object having
$E_\infty^{p, q} = Z_\infty^{p, q}/B_\infty^{p, q}$ in bidegree $(p, q)$,
\item the $p$th graded part $\text{gr}^pH^n(K^\bullet)$ of the
$n$th cohomology object of $K^\bullet$ is a subquotient of
$E_\infty^{p, n - p}$.
\end{enumerate}
\end{lemma}

\begin{proof}
The objects $Z_\infty$, $B_\infty$, and the limit
$E_\infty = Z_\infty/B_\infty$ of
Definition \ref{definition-limit-spectral-sequence}
are bigraded objects of $\mathcal{A}$ by our construction of the
spectral sequence in Lemma \ref{lemma-spectral-sequence-filtered-complex}.
Since $Z_r^{p, q}$, resp.\ $B_r^{p, q}$ is the bidegree $(p, q)$
part of $Z_r$, resp.\ $B_r$, if we assume that
$$
Z_\infty^{p, q} = \bigcap\nolimits_r Z_r^{p, q} =
\bigcap\nolimits_r
\frac{F^pK^{p + q} \cap d^{-1}(F^{p + r}K^{p + q + 1}) + F^{p + 1}K^{p + q}}
{F^{p + 1}K^{p + q}}
$$
and
$$
B_\infty^{p, q} = \bigcup\nolimits_r B_r^{p, q} =
\bigcup\nolimits_r
\frac{F^pK^{p + q} \cap d(F^{p - r + 1}K^{p + q - 1}) + F^{p + 1}K^{p + q}}
{F^{p + 1}K^{p + q}}
$$
exist, then $Z_\infty$ and $B_\infty$ exist with bidegree $(p, q)$
parts $Z_\infty^{p, q}$ and $B_\infty^{p, q}$ (follows from an elementary
argument about unions and intersections of bigraded objects). Thus
$$
E_\infty^{p, q} =
\frac{\bigcap_r (F^pK^{p + q} \cap d^{-1}(F^{p + r}K^{p + q + 1})
+ F^{p + 1}K^{p + q})}
{\bigcup_r (F^pK^{p + q} \cap d(F^{p - r + 1}K^{p + q - 1})
+ F^{p + 1}K^{p + q})}.
$$
where the top and the bottom exist. With $n = p + q$ we have
\begin{equation}
\label{equation-on-top-bigraded}
\Ker(d) \cap F^pK^{n} + F^{p + 1}K^{n}
\subset
\bigcap\nolimits_r
\left(
F^pK^{n} \cap d^{-1}(F^{p + r}K^{n + 1}) + F^{p + 1}K^{n}
\right)
\end{equation}
and
\begin{equation}
\label{equation-at-bottom-bigraded}
\bigcup\nolimits_r
\left(
F^pK^{n} \cap d(F^{p - r + 1}K^{n - 1}) + F^{p + 1}K^{n}
\right)
\subset
\Im(d) \cap F^pK^{n} + F^{p + 1}K^{n}.
\end{equation}
Thus a subquotient of $E_\infty^{p, q}$ is
$$
\frac{\Ker(d) \cap F^pK^{n} + F^{p + 1}K^n}
{\Im(d) \cap F^pK^{n} + F^{p + 1}K^{n}} =
\frac{\Ker(d) \cap F^pK^n}{\Im(d) \cap F^pK^n + \Ker(d) \cap F^{p + 1}K^n}
$$
Comparing with (\ref{equation-graded-cohomology}) we conclude.
\end{proof}

\end{document}
